\section{Methodology: One Architecture, Two Specialisations}
\label{sec:methodology}

We describe the architecture at three levels of abstraction:
domain-agnostic (\Cref{sec:methodology-abstract}), specialised for LOB
(\Cref{sec:methodology-lob}), and specialised for DeFi
(\Cref{sec:methodology-defi}).  \Cref{sec:methodology-loss} states the
composite loss used invariantly in both specialisations.

\subsection{Domain-agnostic DA-BiGRU-CNN}
\label{sec:methodology-abstract}

The DA-BiGRU-CNN architecture
(\emph{Domain-Aware Bidirectional GRU with Convolutional fusion}) takes
a fixed-length window of $T$ time steps and produces an
$H$-step-ahead point forecast for one or more targets.  Inputs are
partitioned into two disjoint subspaces $\mathcal{X}_A$ and
$\mathcal{X}_B$, chosen so that each subspace is a domain-natural
sub-process and their joint observation is informative for the target.
Each subspace is processed by a separate Bidirectional GRU
\cite{cho2014gru} of hidden dimension $d=96$ stacked two layers deep.
The two branch outputs are concatenated and passed through a multi-scale
1-D convolutional fusion head with parallel kernel sizes
$k \in \{3, 5, 7\}$, then projected through a two-layer MLP
$192 \to 96 \to 48 \to H$ to produce the forecast vector.  The total
parameter budget is $\ModelParams$, well below the depth-scaling regime
explored in \cite{briola2021deep, briola2025microstructure}.  The
dual-branch design is conceptually a financial analogue of the
two-stream networks of \cite{simonyan2014two} for video action
recognition: two complementary informational streams fused at a
late stage, not concatenated at the input.

\begin{figure}[t]
  \centering
  % Placeholder; the dual-panel architecture figure is staged in Task 22.
  \fbox{\parbox{0.92\linewidth}{\centering
    \vspace{8pt}
    \textit{Architecture figure placeholder.}\\
    DA-BiGRU-CNN: branch~A (BiGRU, $d{=}96$, 2 layers) over
    $\mathcal{X}_A$; branch~B (BiGRU, $d{=}96$, 2 layers) over
    $\mathcal{X}_B$; concat $\to$ multi-scale Conv1d ($k\in\{3,5,7\}$)
    $\to$ MLP $192{\to}96{\to}48{\to}H$.
    \vspace{8pt}}}
  \caption{Domain-agnostic schematic of DA-BiGRU-CNN. The same wiring
    is used in both the LOB and DeFi specialisations; only the contents
    of $(\mathcal{X}_A, \mathcal{X}_B)$ and the dimensionality of $H$
    differ. Full dual-panel figure to be substituted in the
    camera-ready.}
  \label{fig:architecture-dual}
\end{figure}

\subsection{Specialisation for LOB (recap)}
\label{sec:methodology-lob}

In the LOB specialisation \cite{solovev2026lob} the subspace pair is
$(\mathcal{X}_A, \mathcal{X}_B) = $ (price-series features,
volume-series features), with an additional embedding $e_t$ of 21
microstructure features (depth imbalance at the top-of-book, queue
lengths, spread, immediacy proxies, tick-aligned signed volume) used to
modulate the fusion stage. The motivation for the pair is that LOB
mid-price is the price-side observable but its short-horizon dynamics
are driven by volume-side imbalance signals --- the two-stream
inductive bias from action-recognition translates directly. The
target is the mid-price change at horizon $H$; the model trains on
sequence length $T=\SeqLen$ ticks.

\subsection{Specialisation for DeFi rates}
\label{sec:methodology-defi}

For DeFi lending the structurally analogous subspace pair is:
\begin{equation}
\mathcal{X}_A = \bigl(\varepsilon_t,\;
                 s_t,\;
                 \tilde{s}_t \bigr),\qquad
\mathcal{X}_B = \bigl(u_t,\;
                  \Delta\textsc{TVL}_t^{24h},\;
                  g_t,\;
                  \cos\phi_t,\;\sin\phi_t \bigr),
\label{eq:defi-pair}
\end{equation}
where, for each protocol,
$\varepsilon_t = r_t^{\text{ann}} - f_{\text{kink}}(u_t)$ is the
\emph{rate residual} (with both terms in annualised scale),
$s_t = r_t^{\text{Aave}} - r_t^{\text{Compound}}$ is the
cross-protocol spread, $\tilde{s}_t$ is its de-trended residual,
$u_t$ is utilization, $\Delta\textsc{TVL}_t^{24h}$ is the 24-hour
percent change in total value locked, $g_t$ is Ethereum gas price
(gwei), and $(\cos\phi_t, \sin\phi_t)$ is the cyclical hour-of-day
encoding.  Sequence length is $T=\SeqLen$ hours; the forecast horizon
is $H=\Horizon$ hours.

The strongest claim of this specialisation is the choice to put the
\emph{rate residual} $\varepsilon = r - f_{\text{kink}}(u)$ rather than
the raw rate $r$ into branch~A.  The kink function $f_{\text{kink}}$ is
known on-chain at decision time; subtracting it removes the
deterministic dependency on $u$ that branch~B already observes, leaving
$\varepsilon$ to carry only the unexplained variance --- the structural
residuals that \cite{baude2025optimal} predict the protocol-optimal rate
deviates from. Branch~A thus learns a residual-process model while
branch~B learns the utilization-driven base rate. The decomposition is
the DeFi analogue of separating price-series from volume-series in LOB
\cite{solovev2026lob}, with $f_{\text{kink}}(u)$ playing the role of the
``deterministic backbone'' that microstructure noise rides on top of.
Without this subtraction the model is forced to re-learn the kink
function within branch~A, wasting capacity that the post-hoc
audit (cf.\ \Cref{sec:limitations}) showed was the proximate cause of
the catastrophic R$^2$ artifact in an earlier pre-fix run.

\subsection{Composite loss}
\label{sec:methodology-loss}

The training objective is the convex combination
\begin{equation}
\mathcal{L}\bigl(\hat y, y\bigr)
  = \alpha\,\textsc{MSE}(\hat y, y)
  + \beta\,\bigl(1 - \rho_w(\hat y, y)\bigr)
  + \gamma\,Q_{90}(\hat y, y),
\label{eq:composite-loss}
\end{equation}
with $(\alpha, \beta, \gamma) = (0.4, 0.5, 0.1)$ in both domains.  Here
$\rho_w$ is the weighted Pearson correlation, with weights $w_t$
proportional to $|y_t|$ to over-emphasise large moves (the metric used
in LOB evaluation \cite{solovev2026lob} and equally diagnostic in
lending-rate forecasting where small-move noise dominates row counts).
$Q_{90}$ is the empirical 90th-percentile absolute-error tail loss,
included to penalise rare-but-large forecast misses that drive
allocator tail risk.  The weighted-Pearson term, not MSE, is what the
network gets credit for; MSE acts as a regulariser to keep predictions
on the correct scale and $Q_{90}$ adds a tail-aware safety net.  The
key point for cross-domain transfer is that
\Cref{eq:composite-loss} contains no domain-specific terms; the same
formula trains both specialisations.
